To check whether the direction of the observation-subset effect is specific to PPO, a targeted SAC comparison repeats the following pairs of four-dimensional observations: $\mathcal O_{12}$ versus $\mathcal O_{23}$ at Node 2, and $\mathcal O_{13}$ versus $\mathcal O_{23}$ at Node 3. Each configuration is trained independently in three runs at the same nominal budget of $4\times10^5$ transitions. The SAC settings used in the controller comparison below are unchanged apart from the observation input; no additional tuning or checkpoint selection is performed. All 12 final policies are evaluated deterministically on the same 20 initial conditions as the PPO paired-subset comparison. This check uses nominal dynamics only. The subsequent configuration-selection example uses PPO policies, so its candidate set is unchanged.

\begin{table}[pos=htbp,width=\linewidth]
\centering\small
\setlength{\tabcolsep}{3pt}
\caption{SAC four-dimensional observation cross-check. Each value is mean $\pm$ sample SD across three training-run means over 20 common initial conditions, shared with the PPO paired-subset tests in Table~\ref{tab:observation_identity}.}
\label{tab:sac_observation_check}
\begin{tabular}{@{}clcc@{}}
\toprule
Node & Subset & MAE & $E_u$ \\
\midrule
2 & $\mathcal O_{12}$ & $0.4637\pm0.0659$ & $40.73\pm19.95$ \\
2 & $\mathcal O_{23}$ & $0.1350\pm0.0101$ & $7.17\pm0.69$ \\
3 & $\mathcal O_{13}$ & $0.0562\pm0.0067$ & $7.54\pm0.79$ \\
3 & $\mathcal O_{23}$ & $0.0840\pm0.0059$ & $10.57\pm0.54$ \\
\bottomrule
\end{tabular}
\end{table}


Table~\ref{tab:sac_observation_check} reproduces the direction of the PPO mean-error comparison. At Node 2, replacing $\mathcal O_{12}$ with $\mathcal O_{23}$ reduces SAC MAE from 0.4637 to 0.1350; at Node 3, replacing $\mathcal O_{13}$ with $\mathcal O_{23}$ increases it from 0.0562 to 0.0840. For each node, the same error ordering holds in all three SAC training runs. The agreement across PPO and SAC strengthens the evidence for placement-dependent observation preferences in the tested pairs. The SAC comparison focuses on these original and alternative subsets, complementing the complete paired-subset evaluation performed with PPO.

Control effort remains algorithm dependent. The Node 2 substitution reduces SAC energy from 40.73 to 7.17, whereas PPO energy changes little on average (9.60 to 9.64). At Node 3, energy increases for both algorithms, but the SAC increase (7.54 to 10.57) is much smaller than the PPO increase (6.85 to 142.70), which includes substantial training-run variability. Mean error ordering therefore transfers more consistently than the magnitude of the effort penalty. All 240 SAC test trajectories remained within the state-protection boundary.
